% The fusion pilot in full. Input by main.tex inside the appendix. Section 6 states the pilot
% in one page; this holds its theory, derivation, setup and remaining results. Every measured
% value is a macro from numbers.tex, read from results/fusion_experiment.json.
% =========================================================================================
\section{The exploratory fusion pilot in full}
\label{app:fusion}

\subsection{The theory applied, and the corollary}
\label{app:fusiontheory}

When two classifiers see conditionally independent inputs, the Bayes decision under equal
priors multiplies their posteriors, the product rule~\citep{kittler1998combining}. % [LIT] kittler1998combining, Section 2.1
When all that is known of each classifier is its accuracy $a_k$, the optimal rule for a choice
between two alternatives is a weighted vote with weights proportional to
$\log\big(a_k/(1-a_k)\big)$~\citep{shapley1984optimizing}, % [LIT] shapley1984optimizing, Theorem II
a result Shapley and Grofman credit also to Nitzan and Paroush~\citep{nitzan1982optimal}, and % [LIT] shapley1984optimizing, note 8; nitzan1982optimal was not opened
which Chair and Varshney give for sensor decisions weighted ``according to the reliability of
the detector''~\citep{chair1986optimal}. % [LIT] chair1986optimal, abstract only
For $C$ balanced classes with errors spread evenly over the wrong ones the weight is
$\ln\big((C-1)a_k/(1-a_k)\big)$~\citep{kuncheva2014weighted,li2014error}. % [LIT] kuncheva2014weighted, Eq. (15); li2014error, Corollary 9
That expression is the minimum logit width of Equation~\eqref{eq:tmin}. The standard remedy
for sources whose confidences cannot be trusted is to calibrate each one and then
combine~\citep{brummer2007fusion,kumar2022calibrated}. % [LIT] brummer2007fusion, Section IV; kumar2022calibrated, Proposition 4.1

Call a count readout \emph{saturated} on an input when its winning class fired at all $T$
steps and no other class fired. Its output is then the one vector that attains the ceiling of
Equation~\eqref{eq:ceiling}, the same for every saturated input and every network, up to which
class wins. Let $S$ be the event that both inputs saturate, and $a_{k,S}$ the accuracy of
input $k$ on $S$.
\begin{enumerate}
\item[(a)] On $S$ the two reports are permutations of each other. A fusion rule that uses only
  the two reports, treats the inputs alike, treats the classes alike and rates a class
  strictly higher when an input does gives the two winners equal scores. With uniform
  tie-breaking it is an unweighted vote: a coin flip whenever the inputs disagree. The
  product, sum and max rules are of this kind. Ju et al.\ observe the like for networks whose
  confidences have all saturated towards one, where averaging ``is identical to the majority
  vote''~\citep{ju2018relative}. % [LIT] ju2018relative, Section 4.4.4
\item[(b)] Against a rule that sides with input A on every disagreement in $S$, such a rule
  loses the amount in Equation~\eqref{eq:fusioncost}. This is an identity: where the two agree
  they are right or wrong together, so the difference in their accuracies on $S$ is carried by
  the disagreements, and the coin gives up half of it.
\item[(c)] A scale $s_k$ on input $k$'s counts with
  $s_kT=\ln\big((C-1)a_k/(1-a_k)\big)$ makes the product rule on $S$ the reliability-weighted
  vote of the theory above. A scale common to both inputs restores nothing.
\end{enumerate}
Parts (a) to (c) are checked by enumeration and by simulation in the accompanying code.

\subsection{Setup}
\label{app:fusionsetup}

No network was trained. Input A is an N-MNIST test sample read by a trained network; input B
is a different test sample of the same class, read by the same network after each event is
removed with probability \fuRate, so the two are independent given the class. The rate was % src: results/fusion_dropout_rate.json
fixed on the calibration split before any fused prediction was computed, as the grid value
that lowers calibration accuracy by closest to ten points (\fuCalClean{} to \fuCalDropped). % src: results/fusion_experiment.json dropout
Three rules are compared over \fuPairings{} random pairings of the \nmSplitTest{} test samples: % src: results/fusion_experiment.json pairings; results/nmnist_T2_seed0_split.npz
R0, the product rule on the raw softmax outputs; R1, the product rule after temperature scaling
each input on calibration data; and R2, the reliability-weighted vote, its weights from each
input's calibration accuracy. Intervals are 95\% bootstrap intervals over test samples. Four
predictions were recorded before the run (Table~\ref{tab:predictions}).

\subsection{Results in full}
\label{app:fusionresults}

At $T=2$ calibrating each input first gives $\delta$ of \fuTwoDeltaMin{} to \fuTwoDeltaMax{} % src: results/fusion_experiment.json delta
across seeds, with the lower end of every interval at or above \fuTwoCiLowMin. The gain is % src: results/fusion_experiment.json delta_ci
small in absolute terms because both inputs are accurate on N-MNIST and the raw product rule
already beats either alone; in relative terms it takes the fused error from \fuTwoErrZero\% to % src: results/fusion_experiment.json acc_r0, acc_r1
\fuTwoErrOne\%. The saturated pairs contribute \fuCountTwoDeltaOnS{} of $\delta$, % src: results/fusion_experiment.json delta_on_s
\fuTwoShareOnS\% of it, and Equation~\eqref{eq:fusioncost} gives \fuCountTwoPredicted{} from $P(S)$ % src: results/fusion_experiment.json predicted_on_s
and the two accuracies on $S$, \fuTwoAccAOnS{} and \fuTwoAccBOnS. The identity describes R1 % src: results/fusion_experiment.json accuracy_a_on_s, accuracy_b_on_s
only where temperature scaling gives the clean input the larger scale, which it did in all
\fuTwoScaleOrderRuns{} runs at $T=2$ (temperatures \fuTwoTempA{} against \fuTwoTempB). The rest % src: results/fusion_experiment.json r1_gives_a_the_larger_scale, temperature_a, temperature_b
of $\delta$, \fuTwoDeltaOffS, comes from pairs outside $S$, which the corollary does not % src: results/fusion_experiment.json delta_off_s
describe. The fused confidence of the raw product rule has an ECE of \fuTwoEceRZeroW{} with
equal-width and \fuTwoEceRZeroM{} with equal-mass bins, and after per-input calibration % src: results/fusion_experiment.json ece_r0_ew, ece_r0_em
\fuTwoEceROneW{} and \fuTwoEceROneM. % src: results/fusion_experiment.json ece_r1_ew, ece_r1_em

From $T=4$ the mean of $\delta$ is \fuCountFourDeltaFine, \fuCountSixDeltaFine, % src: results/fusion_experiment.json predictions.14 mean_delta_by_t
\fuCountEightDeltaFine{} and \fuCountTenDeltaFine{} at $T=4$, 6, 8 and 10, and the interval % src: results/fusion_experiment.json predictions.14 mean_delta_by_t
includes zero in all \fuHighRuns{} of those runs (Table~\ref{tab:fusionhigh}). Saturation is % src: results/fusion_experiment.json delta_ci
still common at $T=4$, where $P(S)$ is \fuCountFourPS, but two saturated inputs almost never % src: results/fusion_experiment.json p_s
disagree there: the joint probability falls from \fuTwoPSDisagree{} at $T=2$ to % src: results/fusion_experiment.json p_s_and_disagree
\fuFourPSDisagree. Prediction 14 required the mean to fall strictly through $T=2$, 4 and 6 and % src: results/fusion_experiment.json p_s_and_disagree
to be lower still at $T=8$ and 10. The means at $T=6$ and $T=10$ are equal, so it is
\verdictFourteen; its other clause, intervals including zero at $T=8$ and 10, was met. The % src: results/fusion_experiment.json predictions.14
prediction asked for an ordering among values that are indistinguishable from zero. For the
accumulator, which has no cap, $\delta$ lies between \fuAccumDeltaMin{} and \fuAccumDeltaMax{} % src: results/fusion_experiment.json delta, accumulator
in every run, but in \fuAccumExcl{} of \fuAccumRuns{} runs its interval excludes zero, at % src: results/fusion_experiment.json predictions.15
$T=\fuAccumExclT$ with $\delta$ of \fuAccumExclDelta{} (\fuAccumExclLo{} to \fuAccumExclHi): % src: results/fusion_experiment.json delta_ci, accumulator
temperature scaling made that fusion slightly worse. Prediction 15 allowed none, and is
\verdictFifteen. % src: results/fusion_experiment.json predictions.15

Two further observations qualify the picture. The reliability-weighted vote, R2, is not the
better rule overall. With two voters it is input A alone, and at $T=2$ its accuracy is
\fuTwoRZeroMinusRTwo{} below the raw product rule's, because outside $S$ the counts carry graded % src: results/fusion_experiment.json r2_minus_r0
evidence that a vote discards. Part (b) compares rules on $S$ only. And temperature scaling does
not always favour the more accurate input: in all \fuLateScaleB{} of the \fuLateRuns{} count % src: results/fusion_experiment.json r1_gives_a_the_larger_scale
runs at $T\ge6$ it gave the event-dropped input the larger scale, since dropping events lowers
the counts themselves. There Equation~\eqref{eq:fusioncost} has the wrong sign for R1, and it
falls outside the measured interval in \fuPredMissAll{} runs, each resting on fewer than ten % src: results/fusion_experiment.json predicted_within_on_s_interval
saturated disagreements in 200{,}000 pairs. % 10,000 test samples x 20 pairings: results/fusion_experiment.json n_test, pairings

\begin{table}[ht]
\centering
\caption{The fusion pilot on N-MNIST at every $T$, mean over three seeds and \fuPairings{} % src: results/fusion_experiment.json pairings
pairings. A and B are the single-input accuracies; R0, R1 and R2 the fused accuracies of the
three rules; $P(S)$ the share of pairs on which both inputs saturate; ``on $S$'' the part of
$\delta$ from those pairs and ``Eq.~\eqref{eq:fusioncost}'' its predicted value; ``excl.'' the
seeds whose interval on $\delta$ excludes zero. An accumulator has no counts, so nothing
saturates and those columns do not apply (n/a). The rows for $T=2$ repeat Table~\ref{tab:fusion}.}
\label{tab:fusionhigh}
\small
\setlength{\tabcolsep}{4pt}
\begin{tabular}{lcrrrrrrrrrc}
\toprule
readout & $T$ & A & B & R0 & R1 & R2 & $P(S)$ & $\delta$ & on $S$ & Eq.~\eqref{eq:fusioncost} & excl. \\
\midrule
count & 2  & \fuCountTwoAccA   & \fuCountTwoAccB   & \fuCountTwoRZero   & \fuCountTwoROne   & \fuCountTwoRTwo   & \fuCountTwoPS   & \fuCountTwoDelta   & \fuCountTwoDeltaOnS   & \fuCountTwoPredicted   & \fuCountTwoExcl   \\ % src: results/fusion_experiment.json
count & 4  & \fuCountFourAccA  & \fuCountFourAccB  & \fuCountFourRZero  & \fuCountFourROne  & \fuCountFourRTwo  & \fuCountFourPS  & \fuCountFourDelta  & \fuCountFourDeltaOnS  & \fuCountFourPredicted  & \fuCountFourExcl  \\ % src: results/fusion_experiment.json
count & 6  & \fuCountSixAccA   & \fuCountSixAccB   & \fuCountSixRZero   & \fuCountSixROne   & \fuCountSixRTwo   & \fuCountSixPS   & \fuCountSixDelta   & \fuCountSixDeltaOnS   & \fuCountSixPredicted   & \fuCountSixExcl   \\ % src: results/fusion_experiment.json
count & 8  & \fuCountEightAccA & \fuCountEightAccB & \fuCountEightRZero & \fuCountEightROne & \fuCountEightRTwo & \fuCountEightPS & \fuCountEightDelta & \fuCountEightDeltaOnS & \fuCountEightPredicted & \fuCountEightExcl \\ % src: results/fusion_experiment.json
count & 10 & \fuCountTenAccA   & \fuCountTenAccB   & \fuCountTenRZero   & \fuCountTenROne   & \fuCountTenRTwo   & \fuCountTenPS   & \fuCountTenDelta   & \fuCountTenDeltaOnS   & \fuCountTenPredicted   & \fuCountTenExcl   \\ % src: results/fusion_experiment.json
\midrule
accumulator & 2  & \fuAccumTwoAccA & \fuAccumTwoAccB & \fuAccumTwoRZero & \fuAccumTwoROne & \fuAccumTwoRTwo & n/a & \fuAccumTwoDelta & n/a & n/a & \fuAccumTwoExcl \\ % src: results/fusion_experiment.json
accumulator & 6  & \fuAccumSixAccA & \fuAccumSixAccB & \fuAccumSixRZero & \fuAccumSixROne & \fuAccumSixRTwo & n/a & \fuAccumSixDelta & n/a & n/a & \fuAccumSixExcl \\ % src: results/fusion_experiment.json
accumulator & 10 & \fuAccumTenAccA & \fuAccumTenAccB & \fuAccumTenRZero & \fuAccumTenROne & \fuAccumTenRTwo & n/a & \fuAccumTenDelta & n/a & n/a & \fuAccumTenExcl \\ % src: results/fusion_experiment.json
\bottomrule
\end{tabular}
\end{table}

\subsection{What follows for fusion methods}
\label{app:fusionmethods}

Two readings follow, and neither is a measurement. First, a fusion weight computed from an
input's own output is a function of that output, so on $S$ it takes one value. Zhang et al.\
and Cao et al.\ bound the generalisation error of dynamic fusion, and the bound improves on
static fusion's when each modality's weight is negatively correlated with its
loss~\citep{zhang2023provable,cao2024predictive}; % [LIT] zhang2023provable, Theorem 2; cao2024predictive, Theorem 3.1 and Corollary 3.2
Zhang et al.\ note that a constant weight has zero covariance with every loss. Restricted to
$S$, an output-derived weight is such a constant: the former condition is met only with
equality, the strict form of the latter is not met, and whatever advantage remains must come
from pairs outside $S$. Their theorems are not contradicted, and methods that predict weights
from features instead of outputs are not affected.

Second, Ma et al.\ argue that a
nonlinearity which saturates at high firing rates costs little, since uncertainty is small
there and ``there is little to be gained from exact inference''~\citep{ma2006bayesian}; Beck et % [LIT] ma2006bayesian, Discussion, p. 1436
al.\ propose inhibition to avoid saturation and give no cost for it~\citep{beck2008probabilistic}. % [LIT] beck2008probabilistic
Their setting is a population code for a continuous variable, not this readout. For a count
readout Equation~\eqref{eq:fusioncost} names the case in which saturation is not cheap: when
two saturated sources of unequal reliability disagree. Nothing is claimed here about
perception or behaviour.

Calibrating each input restores the weighting between them and no more. Every saturated
input still reports one confidence, whatever its true posterior, which is the grouping loss
that Perez-Lebel et al.\ show calibration does not remove~\citep{perezlebel2023beyond}. % [LIT] perezlebel2023beyond
Part (c) yields the best rule that knows two accuracies, not the Bayes rule.
